\section{传输层}

\subsection{传输层概述}

\begin{mydef}{传输层的地位与功能}{transport-role}
传输层是五层模型中\textbf{第一个提供端到端（End-to-End）通信}的层次。网络层及以下仅提供\textbf{主机到主机}的通信，传输层在此基础上实现\textbf{进程到进程}的通信。

\textbf{核心功能：}
\begin{enumerate}
    \item \textbf{端到端通信}：将应用层报文分割为报文段（TCP）或用户数据报（UDP），交付给接收方的对应进程。
    \item \textbf{复用与分用}：发送方利用传输层将多个应用进程的数据汇聚到一个网络层（复用）；接收方传输层将收到的数据交付给正确的应用进程（分用）——依据\textbf{端口号}。
    \item \textbf{差错控制}：检测并纠正数据在传输过程中的错误（校验和、确认与重传）。
    \item \textbf{流量控制}：协调发送方与接收方的速率。
    \item \textbf{拥塞控制}（TCP 独有）：防止过多数据注入网络导致拥塞。
\end{enumerate}
\end{mydef}

\begin{attention}
\textbf{408 辨析：} 传输层是「端到端」的，网络层是「点到点」的。端到端 = 源进程 $\to$ 目的进程；点到点 = 当前路由器 $\to$ 下一跳路由器。中间路由器最多处理到网络层，不触及传输层。
\end{attention}

\begin{formula}{端口号}{port-numbers}
端口号是 16 位整数（0--65535），用于传输层复用/分用。一个套接字地址通常由 IP 地址和端口号组成，但还需结合传输层协议解释：同一主机可同时存在相同端口号的 TCP 与 UDP 端点；一条 TCP 连接则由源/目的 IP、源/目的端口组成的四元组标识。

\begin{tablebox}{端口号分类}
\centering
\begin{tabularx}{\linewidth}{|p{0.28\linewidth}|c|X|}
\hline
类型 & 范围 & 说明 \\
\hline
知名端口（Well-known）& 0--1023 & 由 IANA 分配给标准服务（HTTP=80, DNS=53, SMTP=25）\\
注册端口（Registered）& 1024--49151 & 供应用程序注册使用 \\
动态/私有端口 & 49152--65535 & 客户端临时使用（Ephemeral Port）\\
\hline
\end{tabularx}
\end{tablebox}

\textbf{408 常考的知名端口：} HTTP（80）、DNS（53）、SMTP（25）、POP3（110）、FTP（21/20）、Telnet（23）、HTTPS（443）。
\end{formula}

\subsection{UDP —— 用户数据报协议}

\begin{mydef}{UDP 的特点}{udp-features}
UDP（User Datagram Protocol）提供\textbf{无连接、不可靠}的数据传输服务。它在 IP 层之上仅增加了\textbf{端口号}和\textbf{校验和}两个功能。

\begin{itemize}
    \item \textbf{无连接}：发送前无需建立连接，直接发送数据报。
    \item \textbf{不保证交付}：无确认、无重传，数据报可能丢失、重复或失序。
    \item \textbf{面向报文}：发送方 UDP 对应用层交下来的报文，添加首部后直接交付 IP 层。既不合并也不拆分，一次交付一个完整报文（报文边界保留）。
    \item \textbf{无拥塞控制}：发送速率不受网络拥塞限制（适合实时应用）。
    \item \textbf{首部开销小}：仅 8 字节。
\end{itemize}
\end{mydef}

\begin{attention}
\textbf{408 辨析：UDP 面向报文 vs TCP 面向字节流。} UDP 保留报文边界——接收方一次 recvfrom 调用收到一个完整的 UDP 数据报。TCP 不保留报文边界——发送方 3 次 write 各 100 字节，接收方可能一次 read 收到 300 字节，也可能分 3 次各 100 字节。这一区别在选择题中反复出现。
\end{attention}

\begin{conclusion}{UDP 首部格式（8 字节）}{udp-header}
\[
\begin{array}{|c|c|}
\hline
\text{源端口号（16 bit）} & \text{目的端口号（16 bit）} \\
\hline
\text{UDP 长度（16 bit）} & \text{UDP 校验和（16 bit）} \\
\hline
\end{array}
\]

\textbf{字段说明：}
\begin{itemize}
    \item \textbf{源端口号}：发送进程的端口号。若无需回复，可为 0。
    \item \textbf{目的端口号}：接收进程的端口号（必需）。
    \item \textbf{长度}：UDP 数据报的总长度（首部 8B + 数据），单位字节。最小值 8（即无数据）。
    \item \textbf{校验和}：差错检测，覆盖 UDP 首部 + 数据 + \textbf{伪首部}（来自 IP 层的源/目的 IP 地址、协议号、UDP 长度）。
\end{itemize}
\end{conclusion}

\begin{formula}{UDP 校验和的计算}{udp-checksum}
\textbf{UDP 校验和}使用\textbf{伪首部}（Pseudo-header）参与计算：

\begin{tablebox}{UDP 伪首部（12 字节）}
\centering
\begin{tabular}{|l|p{5cm}|}
\hline
字段 & 说明 \\
\hline
源 IP 地址（4B）& 来自 IP 首部 \\
目的 IP 地址（4B）& 来自 IP 首部 \\
全 0（1B）& 填充 \\
协议号（1B）& UDP = 17（TCP = 6）\\
UDP 长度（2B）& 与 UDP 首部中的长度字段相同 \\
\hline
\end{tabular}
\end{tablebox}

\textbf{计算步骤：}
\begin{enumerate}
    \item 构造伪首部 + UDP 首部 + UDP 数据的临时数据块。
    \item 若总长度为奇数字节，末尾补 0 字节。
    \item 按 16 位字进行\textbf{二进制反码求和}（将所有 16 位字以反码加法累加，结果取反码即为校验和）。
    \item 校验和为 0 时填入全 1（$\texttt{0xFFFF}$）——反码表示中全 0 与全 1 等价。
\end{enumerate}
\end{formula}

\subsection{TCP —— 传输控制协议}

\begin{mydef}{TCP 的特点}{tcp-features}
TCP（Transmission Control Protocol）提供\textbf{面向连接的、可靠的、全双工}的字节流传输服务。

\begin{itemize}
    \item \textbf{面向连接}：通信前须经三次握手建立连接，通信后四次挥手释放。
    \item \textbf{可靠传输}：序号 + 确认 + 超时重传 + 校验和，保证数据无差错、不丢失、不重复、按序到达。
    \item \textbf{全双工}：同一连接上双方可同时收发数据。
    \item \textbf{面向字节流}：应用层交付的数据被 TCP 视为无结构的字节流（不保留报文边界）。
    \item \textbf{流量控制}：滑动窗口机制，防止发送方淹没接收方。
    \item \textbf{拥塞控制}：防止过多数据注入网络引起拥塞。
\end{itemize}
\end{mydef}

\subsection{TCP 可靠传输机制}

\begin{conclusion}{序号与确认号}{tcp-seq-ack}
\textbf{序号（Sequence Number）：} TCP 将数据流中的每个字节赋予一个序号。TCP 报文段首部中的「序号」字段是该报文段所携带\textbf{数据部分的第一个字节的序号}。

\textbf{确认号（Acknowledgment Number）：} 期望收到对方\textbf{下一个报文段的第一个数据字节的序号}，即「$\text{已收到连续数据的最后一个字节序号} + 1$」。确认号 $= N$ 表示序号 $N-1$ 及之前的所有数据均已正确收到。

\textbf{累计确认（Cumulative ACK）：} TCP 采用累计确认——确认号 $N$ 表示「$N$ 之前的所有数据已收到」，而非仅确认最后到达的报文段。若报文段乱序到达，接收方仍发送期望的连续序号。

\textbf{初始序号（ISN）：} 连接建立时双方随机选择初始序号（三次握手中交换）。
\end{conclusion}

\begin{attention}
\textbf{408 高频计算：已知发送序号和确认号，推断数据长度。}
\begin{itemize}
    \item 若 A 发送序号 = 101，数据长度 = 200 字节，则下一报文段序号 = 301。
    \item 若 B 发回确认号 = 301，则表示 B 已正确收到序号 $\leqslant$ 300 的所有数据。
    \item \textbf{注意：} TCP 首部中的确认号字段仅在 ACK 标志 = 1 时有效。
\end{itemize}
\end{attention}

\begin{conclusion}{超时重传与 RTT 估计}{tcp-retrans}
TCP 通常为最早的未确认数据维护一个\textbf{重传定时器}：发送含数据的报文段时，若定时器尚未运行则启动；收到推进确认后重启，所有未确认数据均被确认后停止。定时器超时则重传最早的未确认报文段，并按规则退避 RTO。教材中“每个报文段设置定时器”是对这一机制的简化描述，不应理解为实现中为每个报文段各维护一个独立定时器。

\textbf{往返时间 RTT 的估计（Jacobson/Karels 算法）：}
为避免公式过长，记 $E$ 为 EstimatedRTT，$D$ 为 DevRTT，$S$ 为本次 SampleRTT，则
\begin{align}
E_{\mathrm{new}}&=(1-\alpha)E_{\mathrm{old}}+\alpha S,
&\alpha&=0.125,\\
D_{\mathrm{new}}&=(1-\beta)D_{\mathrm{old}}+\beta\lvert S-E_{\mathrm{old}}\rvert,
&\beta&=0.25,\\
\mathrm{RTO}&=E_{\mathrm{new}}+4D_{\mathrm{new}}.&&
\end{align}

\textbf{Karn 算法（重点）：} 当报文段发生重传时，其确认无法区分对应原始发送还是重传副本，因此\textbf{不使用这类有歧义的 SampleRTT 更新 EstimatedRTT}。发生重传超时时，对 RTO 采用指数退避；收到能够取得无歧义 RTT 样本的新数据确认后，再恢复按估计公式更新。

\textbf{408 要点：} 超时重传 + Karn 算法结合考察。选择题经常给出 SampleRTT 序列，其中某一跳是重传的，问更新后的 EstimatedRTT。
\end{conclusion}

\begin{conclusion}{流量控制 —— 滑动窗口}{tcp-flow-control}
TCP 的流量控制通过\textbf{接收窗口（rwnd, Receiver Window）}实现：接收方在 TCP 首部「窗口大小」字段告知发送方自己还有多少可用的接收缓冲区空间。

发送方的发送窗口：
\[
\text{发送窗口} = \min(\text{rwnd}, \text{cwnd})
\]
其中 rwnd 由接收方告知，cwnd 由拥塞控制算法确定。

\textbf{滑动窗口的三种边界指针：}
\begin{center}
\begin{tabular}{|c|c|c|c|c|}
\hline
\multicolumn{2}{|c|}{\textcolor{blue}{已发送已确认}} &
\multicolumn{1}{c|}{\textcolor{red}{已发送未确认}} &
\multicolumn{1}{c|}{\textcolor{green}{可发送（窗口内）}} &
\multicolumn{1}{c|}{不可发送} \\
\hline
\multicolumn{2}{c}{$\longleftarrow\;$} &
\multicolumn{2}{c}{$\longleftarrow\;\;$ 发送窗口 $\;\;\longrightarrow$} &
\\
\hline
\end{tabular}
\end{center}

\begin{itemize}
    \item \textbf{前沿（右边界）}：收到新的 rwnd 时右移（窗口扩大）或被拥塞控制限制。
    \item \textbf{后沿（左边界）}：收到累计确认时右移（窗口滑动）。
    \item \textbf{零窗口}：rwnd = 0 时发送方暂停发送，但会周期性地发送\textbf{零窗口探测报文（Zero-Window Probe）}。
    \item \textbf{糊涂窗口综合征（Silly Window Syndrome）}：发送方每次只发极少数据，效率极低。
    解决方案：Nagle 算法（发送方攒够 MSS 或收到 ACK 才发）、Clark 算法（接收方攒够一定空间才通告窗口）。
\end{itemize}
\end{conclusion}

\subsection{TCP 拥塞控制（408 重中之重）}

\begin{attention}
\textbf{本小节是 408 计算机网络的高频考点。} 常见考法包括画 cwnd 随轮次变化图、填写 cwnd 与 ssthresh，以及区分超时和 3 个重复 ACK 后的处理。作答时必须先采用题目明确给出的 TCP 版本和取整规则。
\end{attention}

\begin{mydef}{拥塞控制的基本概念}{congestion-basics}
拥塞（Congestion）：网络中过多分组导致路由器队列溢出、时延剧增、吞吐量下降的现象。拥塞控制是发送方对注入网络的数据速率进行的自我约束。

\textbf{关键变量：}
\begin{itemize}
    \item \textbf{cwnd（Congestion Window，拥塞窗口）}：发送方根据网络拥塞程度设置的窗口。实际发送窗口 = $\min(\text{rwnd}, \text{cwnd})$。
    \item \textbf{ssthresh（Slow Start Threshold，慢开始门限）}：cwnd 增长方式的分水岭。cwnd $<$ ssthresh 时执行慢开始；cwnd $\geqslant$ ssthresh 时执行拥塞避免。
    \item \textbf{MSS（Maximum Segment Size）}：最大报文段长度（数据部分）。cwnd 和 ssthresh 通常以 MSS 或字节为单位，408 题目中以 MSS 为单位更常见。
\end{itemize}
\end{mydef}

\begin{conclusion}{四种拥塞控制算法详解}{tcp-congestion-algo}
\textbf{TCP 拥塞控制的四种算法形成一个完整的状态机。各状态之间的转换关系如下（文字描述，考试用文字或表格作答即可）：}

\begin{itemize}
    \item 连接建立/超时 $\to$ \textbf{慢开始}（cwnd=1，指数增长）
    \item 慢开始中 cwnd $\geqslant$ ssthresh $\to$ \textbf{拥塞避免}（线性增长）
    \item 拥塞避免中收到 3 个重复 ACK $\to$ \textbf{快重传}；后续窗口变化取决于题设采用的教材化模型或 TCP Reno 快恢复
    \item 快恢复中收到新数据 ACK $\to$ \textbf{拥塞避免}（cwnd=ssthresh，线性增长）
    \item 任何阶段发生\textbf{超时} $\to$ \textbf{慢开始}（ssthresh=cwnd/2, cwnd=1）
\end{itemize}

\textbf{算法一：慢开始（Slow Start）}
\begin{itemize}
    \item \textbf{触发条件：} 连接建立时（cwnd 初始 = 1 MSS），或发生\textbf{超时}事件后重置。
    \item \textbf{cwnd 增长策略：} 每收到一个 ACK，cwnd += 1 MSS（即每个 RTT 后 cwnd 翻倍，指数增长）。
    \item \textbf{退出条件：} cwnd 达到 ssthresh 时转入拥塞避免；或发生超时。
    \item \textbf{超时时的动作：} ssthresh = cwnd/2，cwnd = 1 MSS，重新慢开始。
\end{itemize}

\textbf{算法二：拥塞避免（Congestion Avoidance）}
\begin{itemize}
    \item \textbf{触发条件：} cwnd $\geqslant$ ssthresh。
    \item \textbf{cwnd 增长策略：} 每个 RTT 后 cwnd += 1 MSS（线性增长——加法增大，Additive Increase）。
    \item \textbf{超时时的动作：} ssthresh = cwnd/2，cwnd = 1 MSS，转入慢开始。
\end{itemize}

\textbf{算法三：快重传（Fast Retransmit）}
\begin{itemize}
    \item \textbf{触发条件：} 发送方连续收到 3 个重复 ACK（Duplicate ACK）。
    \item \textbf{动作：} \textbf{立即重传}缺失的报文段（不等超时）。比超时重传快得多。
    \item \textbf{注意：} 收到重复 ACK 说明网络仍有数据在流动，因此不需要像超时那样把 cwnd 降为 1。
\end{itemize}

\textbf{算法四：快恢复（Fast Recovery）}
\begin{itemize}
    \item \textbf{触发条件：} 快重传之后执行。
    \item \textbf{TCP Reno 的动作：} ssthresh 取丢包时飞行中数据量的一半；收到第 3 个重复 ACK 时把 cwnd 暂时置为 ssthresh $+3$ MSS，额外的重复 ACK 可继续使 cwnd 膨胀。收到确认新数据的 ACK 后，将 cwnd 收缩到 ssthresh 并转入拥塞避免。
    \item \textbf{408 教材化计算规则：}不少教材和试题直接把 3 个重复 ACK 后的 cwnd 记为 ssthresh，并从该值进入拥塞避免，省略 Reno 快恢复期间的临时窗口膨胀。若题目未说明，本章习题统一使用这一简化规则；若题目明确指定 Reno，则按上一条计算。
    \item \textbf{若快恢复期间发生超时：} 与普通超时处理相同——ssthresh = cwnd/2，cwnd = 1 MSS，转入慢开始。
\end{itemize}
\end{conclusion}

\begin{conclusion}{拥塞窗口变化示例（408 必须掌握）}{cwnd-graph}
下表展示一个典型 TCP 连接中 cwnd 随传输轮次的变化过程。假定初始 ssthresh = 16 MSS，MSS = 1 KB。

\setlength{\tabcolsep}{3pt}
\begin{tablebox}{cwnd 变化过程示例（初始 ssthresh = 16 MSS）}
\centering
\small
\begin{tabular}{|c|c|c|l|}
\hline
轮次 & cwnd & ssthresh & 状态与事件 \\
\hline
1 & 1 & 16 & 慢开始 \\
2 & 2 & 16 & 慢开始 \\
3 & 4 & 16 & 慢开始 \\
4 & 8 & 16 & 慢开始 \\
5 & 16 & 16 & 慢开始，cwnd$\geqslant$ssthresh，下一轮切换 \\
6 & 17 & 16 & 拥塞避免（线性+1）\\
7 & 18 & 16 & 拥塞避免 \\
\dots & \dots & \dots & \dots \\
— & 24 & 16 & 发生超时！\\
— & 1 & 12 & ssthresh=24/2=12, cwnd=1, 慢开始 \\
— & 2 & 12 & 慢开始 \\
— & 4 & 12 & 慢开始 \\
— & 8 & 12 & 慢开始 \\
— & 12 & 12 & 慢开始，cwnd$\geqslant$ssthresh，切换 \\
— & 13 & 12 & 拥塞避免 \\
— & 14 & 12 & 拥塞避免 \\
— & 15 & 12 & 拥塞避免 \\
— & 16 & 12 & 收到 3 个重复 ACK（快重传）\\
— & 8 & 8 & 教材化规则：ssthresh=16/2=8，cwnd=8 \\
\hline
\end{tabular}
\end{tablebox}

\textbf{408 画图/填表的统一规则：}
\begin{enumerate}
    \item \textbf{超时：} ssthresh = $\max(\text{cwnd}/2, 2\text{ MSS})$，cwnd = 1 MSS，重新慢开始。
    \item \textbf{3 个重复 ACK：}本章习题采用教材化规则，ssthresh 取 cwnd 的一半（整数 MSS 的取整方式由题设说明），cwnd 置为 ssthresh 后进入拥塞避免。只有题目明确指定 TCP Reno 快恢复时，才在第 3 个重复 ACK 到来时使用 ssthresh$+3$ 的临时膨胀窗口。
    \item cwnd 增长：慢开始阶段每个 RTT $\times 2$（指数）；拥塞避免阶段每个 RTT $+1$（线性）。
\end{enumerate}
\end{conclusion}

\begin{attention}
\textbf{408 拥塞控制的高频考法：}
\begin{itemize}
    \item \textbf{画 cwnd 变化曲线：} 给定一组事件序列（收到 ACK / 超时 / 3 dup ACK），画出 cwnd 随时间/轮次的变化。
    \item \textbf{填表：} 给定表格，逐轮次填写 cwnd 和 ssthresh。
    \item \textbf{计算吞吐量：} 给定 cwnd 曲线和 RTT，计算平均吞吐量。
    \item \textbf{区分超时后的处理和快重传后的处理：} 前者 cwnd 降为 1（重新慢开始），后者 cwnd 降至 ssthresh（快恢复）。
    \item \textbf{发送窗口 = $\min(\text{cwnd}, \text{rwnd})$：} 实际发送窗口由拥塞控制和流量控制共同决定。
\end{itemize}
\end{attention}

\subsection{TCP 连接管理}

\begin{conclusion}{三次握手（Three-Way Handshake）}{tcp-handshake}
TCP 连接的建立需要三次报文交换，称为「三次握手」。

\vspace{4pt}
\begin{center}
\begin{tabular}{lcl}
\textbf{客户端} & & \textbf{服务器} \\
\hline
CLOSED & & LISTEN \\
$\downarrow$ & & \\
SYN-SENT & $\xrightarrow{\text{SYN}=1,\;\text{seq}=x}$ & \\
& & SYN-RCVD \\
& $\xleftarrow{\text{SYN}=1,\;\text{ACK}=1,\;\text{seq}=y,\;\text{ack}=x+1}$ & \\
ESTABLISHED & & \\
& $\xrightarrow{\text{ACK}=1,\;\text{seq}=x+1,\;\text{ack}=y+1}$ & \\
& & ESTABLISHED \\
\end{tabular}
\end{center}
\vspace{4pt}

\textbf{三步详解：}
\begin{enumerate}
    \item \textbf{客户端 $\to$ 服务器：SYN = 1, seq = $x$（ISN\_C）}。客户端进入 \texttt{SYN-SENT} 状态。
    \item \textbf{服务器 $\to$ 客户端：SYN = 1, ACK = 1, seq = $y$（ISN\_S）, ack = $x+1$}。服务器进入 \texttt{SYN-RCVD} 状态。
    \item \textbf{客户端 $\to$ 服务器：ACK = 1, seq = $x+1$, ack = $y+1$}。双方进入 \texttt{ESTABLISHED} 状态。
\end{enumerate}

\textbf{前两次握手消耗两个序号（$x$ 和 $y$），第三次握手可不消耗序号（若只发 ACK）。}

\textbf{为什么是三次？}
\begin{itemize}
    \item 防止已失效的连接请求报文段突然到达服务器，导致服务器错误地建立连接。
    \item 三次握手确保\textbf{双方都确认了对方的发送能力和接收能力}：客户端在第 3 步确认服务器的收+发，服务器在第 2 步确认客户端的发送能力并在第 3 步确认自己的发送被收到。
\end{itemize}
\end{conclusion}

\begin{conclusion}{四次挥手（Four-Way Wavehand）}{tcp-wavehand}
TCP 连接的释放需要四次报文交换（因为 TCP 是全双工的，两个方向需分别关闭）。

\vspace{4pt}
\begin{center}
\begin{tabular}{lcl}
\textbf{客户端} & & \textbf{服务器} \\
\hline
ESTABLISHED & & ESTABLISHED \\
$\downarrow$ & & \\
FIN-WAIT-1 & $\xrightarrow{\text{FIN}=1,\;\text{seq}=u}$ & \\
& & CLOSE-WAIT \\
& $\xleftarrow{\text{ACK}=1,\;\text{seq}=v,\;\text{ack}=u+1}$ & \\
FIN-WAIT-2 & & \\
& & LAST-ACK \\
& $\xleftarrow{\text{FIN}=1,\;\text{ACK}=1,\;\text{seq}=w,\;\text{ack}=u+1}$ & \\
TIME-WAIT & & \\
& $\xrightarrow{\text{ACK}=1,\;\text{seq}=u+1,\;\text{ack}=w+1}$ & \\
$\downarrow$（等待 2MSL）& & \\
CLOSED & & CLOSED \\
\end{tabular}
\end{center}
\vspace{4pt}

\textbf{四步详解：}
\begin{enumerate}
    \item \textbf{客户端 $\to$ 服务器：FIN = 1, seq = $u$}。客户端进入 \texttt{FIN-WAIT-1}。
    \item \textbf{服务器 $\to$ 客户端：ACK = 1, ack = $u+1$}。服务器进入 \texttt{CLOSE-WAIT}，客户端收到后进入 \texttt{FIN-WAIT-2}。（此时客户端 $\to$ 服务器方向已关闭，服务器仍可发送数据。）
    \item \textbf{服务器 $\to$ 客户端：FIN = 1, ACK = 1, seq = $w$, ack = $u+1$}。服务器进入 \texttt{LAST-ACK}。
    \item \textbf{客户端 $\to$ 服务器：ACK = 1, ack = $w+1$}。客户端进入 \texttt{TIME-WAIT}，等待 \textbf{2MSL}（Maximum Segment Lifetime，报文最大生存时间，通常 2--4 分钟）后进入 \texttt{CLOSED}。服务器收到 ACK 后立即进入 \texttt{CLOSED}。
\end{enumerate}
\end{conclusion}

\begin{attention}
\textbf{TIME-WAIT 等待 2MSL 的原因（408 常考简答）：}
\begin{enumerate}
    \item \textbf{保证最后一个 ACK 能到达服务器：} 若该 ACK 丢失，服务器会重发 FIN，客户端需要在 TIME-WAIT 期间能够重发 ACK。
    \item \textbf{使本连接持续时间内产生的所有报文段都从网络中消失：} 经过 2MSL，本次连接的所有残留报文都会在网络中被丢弃，不会干扰新连接。
\end{enumerate}
\end{attention}

\begin{conclusion}{TCP 有限状态机（完整）}{tcp-fsm}
\small
\begin{tablebox}{TCP 状态转换总表}
\centering
\footnotesize
\begin{tabular}{|l|l|p{4.5cm}|}
\hline
状态 & 说明 & 典型转换 \\
\hline
CLOSED & 无连接（起始/终点）& 被动打开 $\to$ LISTEN；主动打开 $\to$ SYN-SENT \\
LISTEN & 服务器等待连接 & 收到 SYN $\to$ SYN-RCVD \\
SYN-SENT & 客户端已发 SYN，等待匹配 & 收到 SYN+ACK $\to$ ESTABLISHED \\
SYN-RCVD & 服务器已收到 SYN，发回 SYN+ACK & 收到 ACK $\to$ ESTABLISHED \\
ESTABLISHED & 连接建立，正常通信 & 收到 FIN $\to$ CLOSE-WAIT；主动关闭 $\to$ FIN-WAIT-1 \\
FIN-WAIT-1 & 主动关闭方已发 FIN & 收到 ACK $\to$ FIN-WAIT-2；收到 FIN+ACK $\to$ TIME-WAIT \\
FIN-WAIT-2 & 对方已确认本方的 FIN & 收到 FIN $\to$ TIME-WAIT \\
CLOSE-WAIT & 被动关闭方收到 FIN，等待应用关闭 & 发送 FIN $\to$ LAST-ACK \\
LAST-ACK & 被动关闭方已发 FIN，等待 ACK & 收到 ACK $\to$ CLOSED \\
TIME-WAIT & 等待 2MSL & 超时 $\to$ CLOSED \\
\hline
\end{tabular}
\end{tablebox}
\end{conclusion}

\subsection{TCP 首部格式}

\begin{conclusion}{TCP 报文段首部（$\geqslant$ 20 字节）}{tcp-header-format}
{\small
\[
\begin{array}{|c|c|c|c|c|}
\hline
\multicolumn{5}{|c|}{\text{源端口号（16 bit）}\hspace{1.5cm}\text{目的端口号（16 bit）}} \\
\hline
\multicolumn{5}{|c|}{\text{序号 Sequence Number（32 bit）}} \\
\hline
\multicolumn{5}{|c|}{\text{确认号 Acknowledgment Number（32 bit）}} \\
\hline
\text{数据偏移(4)} & \text{保留(6)} & \text{控制位(6)} & \multicolumn{2}{c|}{\text{窗口大小（16 bit）}} \\
\hline
\multicolumn{2}{|c|}{\text{校验和（16 bit）}} & \multicolumn{3}{c|}{\text{紧急指针（16 bit）}} \\
\hline
\multicolumn{5}{|c|}{\text{选项（可变长度，最多 40 字节）}} \\
\hline
\end{array}
\]
}

\textbf{各字段说明（408 重点）：}
\begin{itemize}
    \item \textbf{源端口号 / 目的端口号（各 16 bit）：} 标识发送/接收进程。
    \item \textbf{序号（32 bit）：} 本报文段数据部分第一个字节的序号。
    \item \textbf{确认号（32 bit）：} 期望收到的下一个字节的序号，仅在 ACK = 1 时有效。
    \item \textbf{数据偏移（4 bit）：} 即 TCP 首部长度，以 4 字节为单位。最小值 5（$5\times4=20$ 字节），最大值 15（$15\times4=60$ 字节）。
    \item \textbf{保留（6 bit）+ 控制位（6 bit）：}
    \begin{itemize}
        \item \textbf{URG}：紧急指针有效。
        \item \textbf{ACK}：确认号有效。连接建立后的所有报文段（除第一个 SYN）都应置 ACK = 1。
        \item \textbf{PSH}：接收方应尽快将数据交付应用层（Push）。
        \item \textbf{RST}：复位连接。用于拒绝非法报文或异常关闭。
        \item \textbf{SYN}：同步序号，用于连接建立。SYN = 1 时不能携带数据，但消耗一个序号。
        \item \textbf{FIN}：发送方数据发送完毕，请求释放连接。FIN = 1 的报文段可携带数据，消耗一个序号。
    \end{itemize}
    \item \textbf{窗口大小（16 bit）：} 接收窗口 rwnd（从 ACK 序号开始，接收方还能接收的字节数）。最大 65535 字节，使用窗口扩大选项可达更大。
    \item \textbf{校验和（16 bit）：} 同 UDP——伪首部 + TCP 首部 + 数据。
    \item \textbf{紧急指针（16 bit）：} 仅在 URG = 1 时有效，指示紧急数据末尾的偏移量。
\end{itemize}
\end{conclusion}

\begin{attention}
\textbf{408 标志位组合的经典考法：}
\begin{itemize}
    \item \textbf{第一个 SYN：} SYN = 1, ACK = 0。
    \item \textbf{第二个 SYN+ACK：} SYN = 1, ACK = 1。
    \item \textbf{第三次握手 ACK：} SYN = 0, ACK = 1。
    \item \textbf{第一个 FIN：} FIN = 1, ACK = 1（通常也携带对之前数据的确认）。
    \item \textbf{RST 攻击/拒绝：} RST = 1, ACK = 0（异常拒绝，不确认）或 RST = 1, ACK = 1（正常拒绝）。
    \item \textbf{SYN 洪泛攻击：} 攻击者发送大量 SYN，伪造源 IP，使服务器 SYN-RCVD 队列满。
\end{itemize}
\end{attention}

\subsection{408 高频考点速查}

\begin{conclusion}{传输层 — 408 常考点}{cn5-keypoints}
\begin{enumerate}
    \item \textbf{TCP 拥塞控制的 cwnd 变化}：慢开始（指数）$\to$ 拥塞避免（线性）$\to$ 超时（cwnd=1）与 3 个重复 ACK（教材化规则下 cwnd=ssthresh）的区别，是常见的画图或填表考点。
    \item \textbf{三次握手与四次挥手的状态转换}：各报文标志位组合、序号/确认号的值、状态名称。
    \item \textbf{TCP 序号和确认号的计算}：已知发送序号和数据长度求下一报文段序号；已知确认号求已接收数据量。
    \item \textbf{TCP 首部格式}：标志位（SYN/ACK/FIN/RST）的组合、数据偏移字段的含义。
    \item \textbf{UDP 首部（8 字节）}：校验和计算（含伪首部）、面向报文 vs TCP 面向字节流。
    \item \textbf{TIME-WAIT 等待 2MSL 的原因}：简答题经典考点。
    \item \textbf{发送窗口 = $\min(\text{rwnd}, \text{cwnd})$}：流量控制与拥塞控制的协作。
    \item \textbf{Karn 算法}：重传报文不参与 RTT 估计，超时时间指数退避。
\end{enumerate}
\end{conclusion}

\newpage
\section{习题}

\subsection{基础过关题}

\begin{enumerate}

\item \textbf{（真题风格·选择题）} 一个TCP连接总是以1 KB的最大段长发送TCP段，发送方有足够多的数据要发送。当拥塞窗口为16 KB时发生了超时，如果接下来的4个RTT时间内的TCP段的传输都是成功的，那么当第4个RTT时间内发送的所有TCP段都得到确认后，拥塞窗口大小是
\begin{center}
\begin{tabular}{ll}
(A) 7 KB & (B) 8 KB \\[6pt]
(C) 9 KB & (D) 16 KB
\end{tabular}
\end{center}

\ifshowsolutions\par\noindent\textbf{解析：} 超时$\to$阈值$=16/2=8$ KB，cwnd$=1$ KB，进入慢开始。RTT1: cwnd$=2$；RTT2: cwnd$=4$；RTT3: cwnd$=8$（到达阈值，进入拥塞避免）；RTT4: cwnd$=9$ KB。\boxed{\textbf{答案：}(C)}\fi

\vspace{8pt}

\item \textbf{（真题风格·选择题）} 主机甲向主机乙发送一个TCP报文段，序号为100，数据部分长度为200 B。主机乙收到后发回的确认号是
\begin{center}
\begin{tabular}{ll}
(A) 100 & (B) 200 \\[6pt]
(C) 300 & (D) 301
\end{tabular}
\end{center}

\ifshowsolutions\par\noindent\textbf{解析：} 确认号$=$期望收到的下一字节序号$=100+200=300$。\boxed{\textbf{答案：}(C)}\fi

\vspace{8pt}

\item \textbf{（真题风格·选择题）} TCP使用慢开始和拥塞避免算法。设TCP的慢开始门限ssthresh的初始值为8（单位为报文段）。当拥塞窗口上升到12时，网络发生超时。采用408教材化模型：超时后将ssthresh置为当前cwnd的一半、cwnd置为1；每轮开始时若cwnd小于ssthresh，则该轮按慢开始增长，否则按拥塞避免增长。从超时重置后的cwnd开始记录5个窗口值，其依次为
\begin{center}
\begin{tabular}{ll}
(A) 1,2,4,8,9 & (B) 1,2,4,6,8 \\[6pt]
(C) 1,2,4,8,10 & (D) 1,2,4,8,12
\end{tabular}
\end{center}

\ifshowsolutions\par\noindent\textbf{解析：} 超时后ssthresh$=12/2=6$，cwnd$=1$。记录值从1开始：$1\to2\to4$；当某轮开始时cwnd$=4<6$，该轮仍按慢开始近似翻倍为8，并不会人为截断为6。此后cwnd已不小于门限，按拥塞避免增至9。因此序列为1,2,4,8,9。\boxed{\textbf{答案：}(A)}\fi

\vspace{8pt}

\item \textbf{（真题风格·选择题）} TCP连接管理中，连接释放的过程称为"四次挥手"。以下关于四次挥手的说法，正确的是
\begin{center}
\begin{tabular}{ll}
(A) FIN报文和ACK报文不能在同一报文中发送 \\[4pt]
(B) 主动关闭方发送FIN后立即进入CLOSED状态 \\[4pt]
(C) 被动关闭方收到FIN后还需要发送数据和FIN \\[4pt]
(D) TIME\_WAIT状态的持续时间是2MSL
\end{tabular}
\end{center}

\ifshowsolutions\par\noindent\textbf{解析：} 主动关闭方发送FIN后进入FIN-WAIT-1，收到ACK后进入FIN-WAIT-2，收到FIN后发送ACK进入TIME-WAIT（2MSL）。\boxed{\textbf{答案：}(D)}\fi

\vspace{8pt}

\item \textbf{（真题风格·选择题）} 主机甲和主机乙已建立TCP连接，甲始终以MSS=1 KB大小的段向乙发送数据，并一直有数据发送；乙每收到一个数据段都会发出一个接收窗口为10 KB的确认段。若甲在$t$时刻发生超时，拥塞窗口为8 KB，则从$t$时刻起，不再发生超时的情况下，经过10个RTT后，甲的发送窗口是
\begin{center}
\begin{tabular}{ll}
(A) 10 KB & (B) 12 KB \\[6pt]
(C) 14 KB & (D) 15 KB
\end{tabular}
\end{center}

\ifshowsolutions\par\noindent\textbf{解析：} 发送窗口$=\min(\text{cwnd},\text{rwnd})$，且 rwnd 始终为10 KB。超时后ssthresh$=4$ KB、cwnd$=1$ KB。按本章教材化的逐 RTT 规则，10个成功 RTT 后 cwnd 依次增长为 $2,4,5,6,7,8,9,10,11,12$ KB；即使拥塞窗口继续增长，接收窗口仍把实际发送窗口限制为 $\min(12,10)=10$ KB。\boxed{\textbf{答案：}(A)}\fi

\end{enumerate}

\subsection{真题风格与综合训练}

\begin{enumerate}[resume]

\item \textbf{（真题风格·选择题）} 主机甲和乙新建TCP连接，甲向乙发送一个起始序号为200的TCP段后，乙发回的确认段中确认号为800，则甲发送的TCP段携带的数据字节数是
\begin{center}
\begin{tabular}{ll}
(A) 600 & (B) 599 \\[6pt]
(C) 800 & (D) 200
\end{tabular}
\end{center}

\ifshowsolutions\par\noindent\textbf{解析：} 确认号$=$期望的下一字节。确认号800，起始序号200，数据长度$=800-200=600$ B。\boxed{\textbf{答案：}(A)}\fi

\vspace{8pt}

\item \textbf{（真题风格·综合题）} TCP 的拥塞窗口 cwnd 和门限 ssthresh 的初始值分别为 1 和 16（单位为 MSS）。本题采用本章的教材化计算规则：收到 3 个重复 ACK 后，ssthresh 取当时 cwnd 的一半，cwnd 置为 ssthresh 后进入拥塞避免。按时间顺序，连续成功传输使 cwnd 分别变为 $2,4,8,16,17,18,19,20$，随后发生 3 次重复 ACK（快重传），之后又经历 2 个成功 RTT。
\begin{enumerate}
    \item[(1)] 画出cwnd随传输轮次变化的折线图；
    \item[(2)] 标注每次状态转换，并说明 3 个重复 ACK 后本题是否单独建模 Reno 快恢复阶段；
    \item[(3)] 第2次成功传输后cwnd的值是多少？
\end{enumerate}

\ifshowsolutions\par\noindent\textbf{解析：}
(1) 慢开始：$1\to2\to4\to8\to16$（达到 ssthresh）；随后拥塞避免：$17\to18\to19\to20$。
3 次重复 ACK 触发快重传：ssthresh$=20/2=10$，按题设的教材化规则，cwnd$=10$，随后继续拥塞避免：$11\to12$。

(2) cwnd 从 1 增至 16 的阶段为慢开始；从 16 增至 20 的阶段为拥塞避免；3 个重复 ACK 触发快重传并把窗口降到 10，此后按拥塞避免线性增长。题干已经明确采用“cwnd 直接置为 ssthresh”的教材化规则，因此这里不再单独套用 TCP Reno 中暂时令 cwnd 增至 $\mathrm{ssthresh}+3$ MSS 的快恢复细节。
(3) 事件发生后的第 2 个成功 RTT 结束时，cwnd$=12$。\fi

\end{enumerate}

\vspace{8pt}
\noindent\textbf{拓展阅读：}
\begin{itemize}
    \item Kurose \& Ross 第3章——TCP可靠传输和拥塞控制的完整讲解与状态图。
    \item 谢希仁 第5章——国内标准术语体系。
\end{itemize}
